% Fichier généré automatiquement par tools/generate_corrections.py.
% Source : CIN/CIN-02-VitesseAcceleration/03_TT/03_TT.tex
% Statut : complet (5/5 question(s) renseignée(s), 0 reprise(s)).
\begin{corrigebox}[Corrigé extrait de la banque]
\CorrigeQuestion{1}
Le point $C$ a un mouvement quelconque dans le plan $\left(A,\vi{0},\vj{0}\right)$.
\CorrigeQuestion{2}
On  a $\vect{AC}=\lambda(t)\vect{i_0}+\mu(t)\vect{j_0}$ et donc, on a directement 
$\left\{
\begin{array}{l}
x_C(t)= \lambda(t) \\
y_C(t)= \mu(t)\\
z_C(t)= 0\\
\end{array}
\right.
$ dans le repère $\repere{A}{i_0}{j_0}{k_0}$.
\CorrigeQuestion{3}
On a $v=R\dot{\theta}(t)$. Par intégration, $ \theta(t)=\dfrac{v}{R}t$ (avec $\theta(t)=\SI{0}{rad}$ pour $t=\SI{0}{s}$).
\CorrigeQuestion{4}
Exprimons la trajectoire du point $C$ : $\vect{AC}=R\vect{e_r} = R\cos\theta(t) \vi{0}+ R\sin\theta(t)  \vj{0}$. Par identification 
$ \lambda(t) = R\cos\theta(t)$ et $\mu(t) = R\sin\theta(t)$.

Au final, 
$\left\{
\begin{array}{l}
\lambda(t) = R\cos\left( \dfrac{v}{R}t\right)\\
\mu(t) = R\sin\left( \dfrac{v}{R}t\right)\\
\end{array}
\right.
$.
\CorrigeQuestion{5}
\begin{lstlisting}
import numpy as np
import matplotlib.pyplot as plt
import math as m
R = 0.1 # m
v = 0.01 # m.s-1 

# Temps pour faire un tour 
T = 2*m.pi*R/v

les_t = np.linspace(0,T,200)
les_lambda = R*np.cos(v/R*les_t)
les_mu = R*np.sin(v/R*les_t)
plt.grid()
plt.plot(les_t,les_lambda,label="$\\lambda(t)$")
plt.plot(les_t,les_mu,label="$\\mu(t)$")
plt.xlabel("Temps ($s$)")
plt.ylabel("Position ($m$)")
plt.legend()
#plt.show()
plt.savefig("03_TT_01_c.pdf")
plt.cla()

plt.grid()
plt.axis("equal")
plt.plot(les_lambda,les_mu,label="Trajectoire de $C$")
plt.legend()
#plt.show()
plt.savefig("03_TT_02_c.pdf")
\end{lstlisting}

\begin{marginfigure}
\includegraphics[width=\linewidth]{03_TT_01_c}
\end{marginfigure}

\begin{marginfigure}
\includegraphics[width=\linewidth]{03_TT_02_c}
\end{marginfigure}
\end{corrigebox}
